\begin{titlepage}
	\begin{center}
		Министерство образования Республики Беларусь \\
		Белорусский Государственный Университет \\
		Физический факультет
	\end{center}
	
	\vspace{3cm}
	
	\begin{center}
		\textbf{Лабораторная работа №1} \\
		\textbf{Обработка результатов измерений}
	\end{center}
	
	\vspace{4cm}
	
	\begin{flushright}
		Выполнила: \\
		Студентка 1 курса 2 группы \\
		Физического факультета БГУ \\
		Щитникова Анна Игоревна \\[1cm]
		Проверил: \\
		Чертко Наталья Викторовна
	\end{flushright}
	
	\vspace{4cm}
	
	\begin{center}
		Вариант 2 \\
		Минск 2019
	\end{center}
\end{titlepage}

\section*{Цель работы}
Изучить метод среднего квадратичного и метод наименьших квадратов.

\section*{Задание №1. Погрешности прямых и косвенных измерений}
Проведены измерения размеров полого цилиндра: $h$ -- высота, $D$ -- внешний диаметр, $d$ -- внутренний диаметр. Все измерения проведены с помощью штангенциркуля.

\begin{figure}[h!]
	\centering
	% Вставьте сюда имя вашего файла с картинкой, например: \includegraphics[width=0.5\textwidth]{cylinder.png}
	\caption{Схема полого цилиндра}
	\label{fig:cylinder}
\end{figure}

Результаты измерений представлены в Таблице 1.

\begin{table}[h!]
	\centering
	\caption{Результаты измерений размеров цилиндра}
	\begin{tabular}{@{}llll@{}}
		\toprule
		№ измерения & $h$, мм & $D$, мм & $d$, мм \\ \midrule
		1 & 70,00 & 29,90 & 25,05 \\
		2 & 69,95 & 29,95 & 24,95 \\
		3 & 70,00 & 30,10 & 25,05 \\
		4 & 70,05 & 30,15 & 24,90 \\
		5 & 70,20 & 29,95 & 25,05 \\
		6 & -     & 29,95 & 25,00 \\
		7 & -     & -     & 24,90 \\
		8 & -     & -     & 25,00 \\ \midrule
		Среднее $\langle x \rangle$ & 70,04 & 30,00 & 24,99 \\
		$\Delta x_{окр}$ & 0,02 & 0,02 & 0,02 \\
		$\Delta x_{пр}$ & 0,03 & 0,03 & 0,03 \\
		$n$ & 5 & 6 & 8 \\
		$t_{n,P}$ & 2,78 & 2,57 & 2,36 \\
		$\Delta x_{сл}$ & 0,12 & 0,10 & 0,05 \\
		$\Delta x$ & 0,13 & 0,11 & 0,06 \\
		$\varepsilon$, \% & 0,19 & 0,37 & 0,24 \\ \bottomrule
	\end{tabular}
\end{table}

Рассчитаем средние значения размеров полого цилиндра по формуле $\langle x \rangle = \frac{\sum x_i}{n}$:
\begin{align*}
	\langle h \rangle &= \frac{70,00 + 69,95 + 70,00 + 70,05 + 70,20}{5} = 70,04 \text{ мм}, \\
	\langle D \rangle &= \frac{29,90 + 29,95 + 30,10 + 30,15 + 29,95 + 29,95}{6} = 30,00 \text{ мм}, \\
	\langle d \rangle &= \frac{25,05 + 24,95 + 25,05 + 24,90 + 25,05 + 25,00 + 24,90 + 25,00}{8} = 24,99 \text{ мм}.
\end{align*}

Рассчитаем погрешность округления. Возьмем интервал округления, равный половине цены деления штангенциркуля ($h = 0,05$ мм), а доверительную вероятность $P = 0,95$. Тогда $\Delta x_{окр}$ для $h$, $D$ и $d$:
$$ \Delta x_{окр} = \frac{h \cdot P}{2} = \frac{0,05 \cdot 0,95}{2} \approx 0,02 \text{ мм}. $$

Рассчитаем приборную погрешность. Предельная погрешность штангенциркуля $\delta = 0,05$ мм:
$$ \Delta x_{пр} = \frac{2}{3} \delta = \frac{2}{3} \cdot 0,05 \approx 0,03 \text{ мм}. $$

Случайная погрешность рассчитывается по формуле:
$$ \Delta x_{сл} = t_{n,P} \sqrt{\frac{\sum_{i=1}^{n}(\langle x \rangle - x_i)^2}{n(n-1)}}. $$
Подставляя значения, получаем: $\Delta h_{сл} = 0,12$ мм, $\Delta D_{сл} = 0,10$ мм, $\Delta d_{сл} = 0,05$ мм.

Абсолютная полная погрешность прямого измерения:
$$ \Delta x = \sqrt{\Delta x_{пр}^2 + \Delta x_{окр}^2 + \Delta x_{сл}^2}. $$
$$ \Delta h = \sqrt{0,03^2 + 0,02^2 + 0,12^2} = 0,13 \text{ мм}; $$
$$ \Delta D = \sqrt{0,03^2 + 0,02^2 + 0,10^2} = 0,11 \text{ мм}; $$
$$ \Delta d = \sqrt{0,03^2 + 0,02^2 + 0,05^2} = 0,06 \text{ мм}. $$

Относительная погрешность $\varepsilon = \frac{\Delta x}{\langle x \rangle} \cdot 100\%$:
$\varepsilon_h = 0,19\%$, $\varepsilon_D = 0,37\%$, $\varepsilon_d = 0,24\%$.

Рассчитаем объем полого цилиндра:
$$ V = \frac{\pi}{4}(\langle D \rangle^2 - \langle d \rangle^2)\langle h \rangle = \frac{3,14159}{4}(30,00^2 - 24,99^2) \cdot 70,04 \approx 15162 \text{ мм}^3. $$

Погрешность определения объема (косвенное измерение):
$$ \Delta V = \sqrt{\left(\frac{\partial V}{\partial h}\Delta h\right)^2 + \left(\frac{\partial V}{\partial D}\Delta D\right)^2 + \left(\frac{\partial V}{\partial d}\Delta d\right)^2}, $$
где частные производные:
$$ \frac{\partial V}{\partial h} = \frac{\pi}{4}(D^2 - d^2), \quad \frac{\partial V}{\partial D} = \frac{\pi}{2}Dh, \quad \frac{\partial V}{\partial d} = -\frac{\pi}{2}dh. $$
$$ \Delta V = \sqrt{\left(\frac{\pi}{4}(30,00^2 - 24,99^2) \cdot 0,13\right)^2 + \left(\frac{\pi}{2} \cdot 30,00 \cdot 70,04 \cdot 0,11\right)^2 + \left(\frac{\pi}{2} \cdot 24,99 \cdot 70,04 \cdot 0,06\right)^2} \approx 400 \text{ мм}^3. $$

\textbf{Итоговый результат Задания 1:}
\begin{itemize}
	\item $h = 70,04 \pm 0,13$ мм; $\varepsilon_h = 0,19\%$; $P = 0,95$;
	\item $D = 30,00 \pm 0,11$ мм; $\varepsilon_D = 0,37\%$; $P = 0,95$;
	\item $d = 24,99 \pm 0,06$ мм; $\varepsilon_d = 0,24\%$; $P = 0,95$;
	\item $V = 15162 \pm 400$ мм$^3$; $\varepsilon_V = 2,6\%$; $P = 0,95$.
\end{itemize}

\section*{Задание №2. Метод наименьших квадратов (линейная зависимость)}
Проведены измерения зависимости $y = ax + b$. Результаты представлены в Таблице 2.

\begin{table}[h!]
	\centering
	\caption{Данные для метода наименьших квадратов}
	\begin{tabular}{@{}llllll@{}}
		\toprule
		№ & $x$ & $y$ & $x^2$ & $y^2$ & $x \cdot y$ \\ \midrule
		1 & 1  & 45,53 & 1   & 2072,98 & 45,53 \\
		2 & 3  & 42,81 & 9   & 1832,70 & 128,43 \\
		3 & 5  & 38,89 & 25  & 1512,43 & 194,45 \\
		4 & 7  & 35,70 & 49  & 1274,49 & 249,90 \\
		5 & 9  & 34,30 & 81  & 1176,49 & 308,70 \\
		6 & 11 & 29,63 & 121 & 877,94  & 325,93 \\
		7 & 13 & 25,79 & 169 & 665,12  & 335,27 \\
		8 & 15 & 24,24 & 225 & 587,58  & 363,60 \\
		9 & 17 & 22,48 & 289 & 505,35  & 382,16 \\
		10 & 19 & 19,41 & 361 & 376,75  & 368,79 \\ \midrule
		Среднее & 10 & 31,88 & 133 & 1088,18 & 270,28 \\ \bottomrule
	\end{tabular}
\end{table}

Рассчитаем дисперсии и коэффициент ковариации:
$$ S_x^2 = \langle x^2 \rangle - \langle x \rangle^2 = 133 - 10^2 = 33; $$
$$ S_y^2 = \langle y^2 \rangle - \langle y \rangle^2 = 1088,18 - 31,88^2 = 71,98; $$
$$ R_{xy} = \langle xy \rangle - \langle x \rangle \langle y \rangle = 270,28 - 10 \cdot 31,88 = -48,52. $$

Коэффициент корреляции:
$$ r = \frac{R_{xy}}{\sqrt{S_x^2 \cdot S_y^2}} = \frac{-48,52}{\sqrt{33 \cdot 71,98}} \approx -1,0. $$

Коэффициенты линейной зависимости:
$$ a = \frac{R_{xy}}{S_x^2} = \frac{-48,52}{33} = -1,47; $$
$$ b = \langle y \rangle - a\langle x \rangle = 31,88 - (-1,47) \cdot 10 = 46,58. $$

Абсолютные случайные погрешности коэффициентов (при $t_{9, 0.95} \approx 2,26$):
$$ \Delta a = t \sqrt{\frac{1}{n-2} \left( \frac{S_y^2}{S_x^2} - a^2 \right)} = 2,26 \sqrt{\frac{1}{8} \left( \frac{71,98}{33} - (-1,47)^2 \right)} \approx 0,10; $$
$$ \Delta b = \Delta a \sqrt{S_x^2 + \langle x \rangle^2} = 0,10 \sqrt{33 + 10^2} \approx 1,15. $$

\textbf{Итоговый результат Задания 2:}
$a = -1,47 \pm 0,10$; $b = 46,58 \pm 1,15$.

\section*{Задание №3. МНК для нелинейной зависимости}
Исходная зависимость имеет вид $y = C x^n$. Прологарифмировав функцию (по основанию 10), получим линейную зависимость:
$$ \lg y = \lg C + n \lg x \implies Y = A X + B, $$
где $Y = \lg y$, $X = \lg x$, $A = n$, $B = \lg C$.

\begin{table}[h!]
	\centering
	\caption{Данные для нелинейного МНК}
	\begin{tabular}{@{}lllllll@{}}
		\toprule
		№ & $x$ & $y$ & $X = \lg x$ & $Y = \lg y$ & $X^2$ & $X \cdot Y$ \\ \midrule
		1 & 1  & 1,647 & 0,000 & 0,217 & 0,000 & 0,000 \\
		2 & 2  & 3,947 & 0,301 & 0,596 & 0,091 & 0,179 \\
		3 & 3  & 7,121 & 0,477 & 0,853 & 0,228 & 0,407 \\
		4 & 4  & 12,616 & 0,602 & 1,101 & 0,362 & 0,663 \\
		5 & 5  & 16,622 & 0,699 & 1,221 & 0,489 & 0,853 \\
		6 & 6  & 22,654 & 0,778 & 1,355 & 0,605 & 1,054 \\
		7 & 7  & 28,450 & 0,845 & 1,454 & 0,714 & 1,229 \\
		8 & 8  & 33,670 & 0,903 & 1,527 & 0,815 & 1,379 \\
		9 & 9  & 39,897 & 0,954 & 1,601 & 0,910 & 1,527 \\
		10 & 10 & 47,136 & 1,000 & 1,673 & 1,000 & 1,673 \\ \midrule
		Среднее & & & 0,656 & 1,160 & 0,521 & 0,896 \\ \bottomrule
	\end{tabular}
\end{table}

Рассчитаем параметры для линеаризованной зависимости:
$$ S_X^2 = \langle X^2 \rangle - \langle X \rangle^2 = 0,521 - 0,656^2 = 0,091; $$
$$ S_Y^2 = \langle Y^2 \rangle - \langle Y \rangle^2 \approx 0,200; $$
$$ R_{XY} = \langle XY \rangle - \langle X \rangle \langle Y \rangle = 0,896 - 0,656 \cdot 1,160 = 0,135. $$

Коэффициенты $A$ и $B$:
$$ A = \frac{R_{XY}}{S_X^2} = \frac{0,135}{0,091} = 1,49; $$
$$ B = \langle Y \rangle - A\langle X \rangle = 1,160 - 1,49 \cdot 0,656 = 0,18. $$

Находим исходные коэффициенты $C$ и $n$:
$$ n = A = 1,49; $$
$$ C = 10^B = 10^{0,18} = 1,52. $$

Погрешности коэффициентов:
$$ \Delta n = \Delta A = t \sqrt{\frac{1}{n-2} \left( \frac{S_Y^2}{S_X^2} - A^2 \right)} \approx 0,06; $$
$$ \Delta B = \Delta A \sqrt{S_X^2 + \langle X \rangle^2} \approx 0,08; $$
$$ \Delta C = C \cdot \ln 10 \cdot \Delta B = 1,52 \cdot 2,303 \cdot 0,08 \approx 0,15. $$

\textbf{Итоговый результат Задания 3:}
$n = 1,49 \pm 0,06$; $C = 1,52 \pm 0,15$.

\section*{Вывод}
В ходе выполнения лабораторной работы №1 были изучены методы обработки данных, такие как методики расчёта погрешностей при прямых и косвенных измерениях. Был изучен метод наименьших квадратов и его применение как для линейной зависимости, так и для нелинейной зависимости.

При выполнении 1-го задания были вычислены средние значения измеренных величин и их погрешности. А также был вычислен объем тела и его погрешность.

При выполнении 2-го задания использовался метод наименьших квадратов (МНК), с помощью которого были вычислены коэффициенты уравнения $y = ax + b$. Были рассчитаны погрешности косвенных измерений.

При выполнении 3-го задания использовался метод логарифмирования и МНК, что позволило рассчитать коэффициенты как для изначальной зависимости $y = Cx^n$, так и для полученной после логарифмирования зависимости. Были рассчитаны погрешности косвенных измерений. По итогу работы погрешности измерений находились в пределах нормы.